% Full independent review from the 4 September 2026 conversation.
% This is a review document, not a replacement manuscript section.
% Build from repository root with pdflatex; keep auxiliary files outside this directory.
\documentclass[11pt,letterpaper]{article}
\usepackage[margin=0.85in,headheight=14pt,headsep=18pt,footskip=25pt]{geometry}
\usepackage[T1]{fontenc}
\usepackage[utf8]{inputenc}
\usepackage{newpxtext,newpxmath}
\usepackage{microtype}
\usepackage{amsmath,amssymb}
\usepackage{booktabs,longtable,array,calc}
\usepackage{pdflscape}
\usepackage{xcolor}
\usepackage{xurl}
\usepackage{enumitem}
\usepackage{needspace}
\usepackage{titlesec}
\usepackage{fancyhdr}
\usepackage{hyperref}
\usepackage{bookmark}
\definecolor{inkblue}{HTML}{234A65}
\hypersetup{colorlinks=true,linkcolor=inkblue,urlcolor=inkblue,
 pdftitle={Independent review of the simplified OLG theory},
 pdfsubject={Formal audit, welfare arbitration, transition verdict, and author decisions}}
\urlstyle{tt}
\setlength{\parindent}{0pt}
\setlength{\parskip}{0.47em}
\linespread{1.045}
\setlength{\emergencystretch}{2.5em}
\setlength{\LTpre}{0.7em}
\setlength{\LTpost}{0.7em}
\setlist[itemize]{leftmargin=1.5em,itemsep=0.35em,topsep=0.35em,parsep=0pt}
\setlist[enumerate]{leftmargin=1.6em,itemsep=0.35em,topsep=0.35em,parsep=0pt}
\providecommand{\tightlist}{\setlength{\itemsep}{0.25em}\setlength{\parskip}{0pt}}
\titleformat{\section}{\Large\bfseries\color{inkblue}}{\thesection}{0.6em}{}
\titleformat{\subsection}{\large\bfseries}{\thesubsection}{0.6em}{}
\titlespacing*{\section}{0pt}{1.3em}{0.6em}
\titlespacing*{\subsection}{0pt}{1em}{0.45em}
\setcounter{secnumdepth}{1}
\setcounter{tocdepth}{1}
\pagestyle{fancy}
\fancyhf{}
\fancyhead[L]{\small\textsc{Simplified OLG theory}}
\fancyhead[R]{\small Independent review}
\fancyfoot[C]{\small\thepage}
\renewcommand{\headrulewidth}{0.25pt}
\widowpenalty=10000
\clubpenalty=10000
\displaywidowpenalty=10000
\allowdisplaybreaks[1]
\begin{document}
\thispagestyle{empty}
{\small\textsc{Fertility\_Spring26}\hfill 4 September 2026\par}
\vspace{1.0cm}
{\LARGE\bfseries Independent review of the simplified OLG theory\par}
\vspace{0.35cm}
{\large Housing, fertility, and intergenerational allocation\par}
\vspace{0.65cm}
This report preserves the full evaluation from the conversation, including the formal
claim audit, the arbitration of welfare claims A--C, the transition assessment, and
the proposed priority ladder. Recommendations remain proposals for discussion;
they are not author decisions or amendments to the model.

The audit refers to the working-tree versions inspected on 4 September 2026.
Source labels and line numbers are retained, with full paths in the source guide.
The protected manuscript has not been amended.
\vspace{0.35cm}
\tableofcontents
\clearpage
\section{Bottom line}\label{bottom-line}

\textbf{I would sign off on the core household algebra and the numerical finite-horizon
construction. I would not sign off on the present package as a baseline misallocation
theorem or a proved transition between two positive steady states.}

The repaired model is substantially sound. Rent timing, mortgage cancellation, old-age
budgets, the capped-old-renter reduction, conditional fertility choices, and steady-state
accounting are consistent. Independent optimization of the original household problems
reproduces the analytical solutions. The numerical transition is credible as a
finite-horizon approximation.

The welfare dispute has a clear resolution:

\begin{itemize}
\tightlist
\item
  \textbf{A is valid under an expanded institution:} bridge credit and targeted,
  intertemporal compensation make the proposed transfer feasible.
\item
  \textbf{B is correct:} the same transfer is infeasible for a down-payment-constrained
  buyer under the baseline closing rule. A and B are not contradictory.
\item
  \textbf{C has a valid zero-transaction-cost version}, conditional on a positive mass of
  unconstrained buyers and interior sellers. It identifies a realization-tax fiscal
  externality. Its proposed extension using only average transaction costs is not
  sufficient under the stated common-rebate restriction.
\end{itemize}

Two additional qualifications matter. The rental-cap derivative is a partial derivative
holding the old-age housing commitment fixed; changing the common cap at both ages adds an
omitted resource effect. And the numerical Jacobian checks do not verify the zero-shock
hypotheses of the fixed-\(J\) proposition.

\textbf{For this job-market paper, I recommend retaining the housing-cost/fertility
mechanism and a compact demographic accounting result, with no welfare proposition in the
main text.} Capital-gains taxation should be an optional extension: the quantitative model
does not contain it, the displayed constrained buyers cannot implement C, and the
empirical application concerns falling desired fertility rather than the appreciation
experiment used here.

The three decisions are:

\begin{enumerate}
\def\labelenumi{\arabic{enumi}.}
\tightlist
\item
  Make the analytical section explain the quantitative mechanism, rather than carry a
  separate welfare contribution.
\item
  Keep the transition numerical and explicitly approximate; do not pursue the nonsmooth
  infinite-horizon theorem for this paper.
\item
  Remove capital-gains taxation from the main analytical benchmark unless its separate
  contribution is important enough to justify its institutional and expositional cost.
\end{enumerate}

At the time of this review, no project files had been changed. The author-controlled
analytical section and appendix remained placeholders.

\clearpage
\section{Model map}\label{model-map}

Here is the timing implied by the actual budgets.

\begingroup\small\setlength{\tabcolsep}{5pt}\renewcommand{\arraystretch}{1.17}

\begin{longtable}[]{@{}
  >{\raggedright\arraybackslash}p{(\columnwidth - 4\tabcolsep) * \real{0.1800}}
  >{\raggedright\arraybackslash}p{(\columnwidth - 4\tabcolsep) * \real{0.4100}}
  >{\raggedright\arraybackslash}p{(\columnwidth - 4\tabcolsep) * \real{0.4100}}@{}}
\toprule\noalign{}
\begin{minipage}[b]{\linewidth}\raggedright
Object
\end{minipage} & \begin{minipage}[b]{\linewidth}\raggedright
Young at \(t\)
\end{minipage} & \begin{minipage}[b]{\linewidth}\raggedright
Old at \(t+1\)
\end{minipage} \\
\midrule\noalign{}
\endhead
\bottomrule\noalign{}
\endlastfoot
Cohorts & Mass \(Y_t\); chooses completed fertility \(n_t\) & This cohort becomes
\(O_{t+1}=Y_t\) \\
Income and wealth & Receives \(y_i\); enters with liquid wealth \(b_i\) & No labor income;
receives matured financial assets net of mortgage repayment \\
Closing & Occurs before income and rebates become available & Existing ownership and
acquisition basis are predetermined \\
Purchase and mortgage & Buys \(h\) for \(P_th\), borrows \(\phi P_th\), posts
\((1-\phi)P_th\leq b_i\) & Mortgage repayment is \(R_f\phi P_th\) \\
Saving & Chooses \(a'\geq0\) & Receives \(R_fa'\); owner net financial wealth is
\(R_fa'-R_f\phi P_th\) \\
Rent & \(r_t h\) is paid at the period's end; its current value is \(qr_th\) & An old
renter similarly pays present-value rent \(qr_{t+1}h^2\) \\
Property tax & Owner of record owes \(\tau^pP_th\) at period end & Pays tax on retained
housing during old age \\
Housing services & Occupies \(h\); children require \(\kappa n\), leaving \(s=h-\kappa n\)
& Renter chooses housing; owner retains \(h^2\leq H\) \\
Living-owner sale & Not an additional choice within young age & Sells \(H-h^2\), paying
\(\ell_{t+1}=\tau^g(P_{t+1}-P_t)_+\) per unit \\
Retention and death & --- & Retained housing supplies services, then is sold by the estate
after its basis resets to the death-date price \\
Estate & Anticipates warm-glow estate utility & Estate \(e\) includes retained housing's
sale value; \(e\geq P_{t+2}h^2\) prevents negative nonhousing estate wealth \\
Government rebate & Receives the current value \(T_t\), unavailable at closing & Receives
\(T_{t+1}\) while alive \\
Descendants & \(Y_{t+1}=\nu\bar n_tY_t\) & Estates do \textbf{not} finance descendants'
entry wealth; new entrants draw fresh \((y,b)\) from \(F\) \\
\end{longtable}

\endgroup

The economy is \textbf{demographically closed but financially open}. Housing clears
domestically; the consumption good and bond trade against an external sector at the fixed
bond price \(q\). Consequently, there is no missing domestic bond-clearing equation. There
would be one if this were described as a closed goods-and-asset economy.

The rental intermediary's primitive cash flow gives:

\[
-P_t+q\bigl(r_t-\tau^pP_t+P_{t+1}\bigr)=0,
\]

and therefore

\[
qr_t=P_t+q\tau^pP_t-qP_{t+1}.
\]

This verifies the household timing. At constant prices, \(r/P=R_f-1+\tau^p\). Assigning
legal property-tax payment to the titleholder does not establish economic incidence; the
substantive asymmetries are realization taxation, exemption of intermediaries, stepped-up
basis, and financing restrictions.

\Needspace{6\baselineskip}
For old owners, define their effective resources and retention cost by:

\[
X_t=a+(P_t-\ell_t)H+T_t,
\qquad
\rho_t=qr_t-\ell_t.
\]

Their budget follows by adding after-tax sales, subtracting property tax on retained
housing, and including its eventual estate value. It counts title and services correctly.

The essential equilibrium objects are prices \(P_t\), rebates \(T_t\), cohort masses, and
the old-state distribution \(G_t\). The state

\[
S_t=(Y_t,O_t,G_t,P_{t-1})
\]

is sufficient because all living old owners bought at the same date. Their ownership taste
need not remain in the state: it affects the original tenure choice but not their old-age
problem.

Two institutional assumptions deserve more prominence:

\begin{itemize}
\tightlist
\item
  Rental limits restrict household housing consumption, including the ability to combine
  rental units. Otherwise the cap could be circumvented.
\item
  Owners cannot separate title from occupancy by renting out retained housing. This
  restriction is essential to the reallocation argument.
\end{itemize}

\clearpage\begin{landscape}

\section{Formal claim audit}\label{formal-claim-audit}

References below point to the working-tree sources. ``Conditional'' means \textbf{correct
only under the stated conditions}, not an unconditional result about the model class.

\begingroup\small\setlength{\tabcolsep}{5pt}\renewcommand{\arraystretch}{1.17}

\begin{longtable}[]{@{}
  >{\raggedright\arraybackslash}p{(\columnwidth - 8\tabcolsep) * \real{0.2050}}
  >{\raggedright\arraybackslash}p{(\columnwidth - 8\tabcolsep) * \real{0.1800}}
  >{\raggedright\arraybackslash}p{(\columnwidth - 8\tabcolsep) * \real{0.2650}}
  >{\raggedright\arraybackslash}p{(\columnwidth - 8\tabcolsep) * \real{0.1150}}
  >{\raggedright\arraybackslash}p{(\columnwidth - 8\tabcolsep) * \real{0.2350}}@{}}
\toprule\noalign{}
\begin{minipage}[b]{\linewidth}\raggedright
Claim
\end{minipage} & \begin{minipage}[b]{\linewidth}\raggedright
Verdict
\end{minipage} & \begin{minipage}[b]{\linewidth}\raggedright
Missing or essential assumptions
\end{minipage} & \begin{minipage}[b]{\linewidth}\raggedright
Evidence
\end{minipage} & \begin{minipage}[b]{\linewidth}\raggedright
Paper consequence
\end{minipage} \\
\midrule\noalign{}
\endhead
\bottomrule\noalign{}
\endlastfoot
Usable bundles \(x=c-\chi n,\ s=h-\kappa n\) & Proved as written & Positive bundles and
positive preference weights & \hyperlink{source-main}{Main:23} & Coherent goods and space
requirements; no double counting \\
Rental zero profit and stationary user cost & Proved as written & End-of-period rent/tax;
untaxed intermediary resale; fixed \(q\) & \hyperlink{source-main}{Main:49} & Retain \\
Young budgets and mortgage cancellation & Conditional & Same borrowing/lending return; no
default; explicit closing chronology & \hyperlink{source-main}{Main:89},
\hyperlink{source-appendix}{Appendix:178} & Down payment is a liquidity restriction, not
an extra lifetime resource cost \\
Old-owner budget and estate floor & Conditional & Estate sells retained title; stepped-up
basis; nonnegative financial estate & \hyperlink{source-main}{Main:115} & Explain the
primitive cash flow before the compressed budget \\
Lemma 1: old-age allocations & Conditional & Positive resources/user costs; owner
retention and estate inequalities; renter branch consistency &
\hyperlink{source-appendix}{Appendix:23} & Correct on the displayed branches \\
Remaining owner corners and zero saving & Incomplete as a solution catalogue & Their KKT
systems are referenced but not fully displayed or implemented &
\hyperlink{source-appendix}{Appendix:65}, \hyperlink{source-appendix}{Appendix:213} & Do
not describe the driver as a general solver \\
Lifetime reduction & Conditional & Interior saving; correct old-age branch &
\hyperlink{source-appendix}{Appendix:87} & Correct repair: capped old renters use a
different saving weight and resource term \\
Proposition 3: unconstrained conditional policies and saving & Conditional & Positive user
cost/resources; all recovered choices satisfy their branch inequalities &
\hyperlink{source-appendix}{Appendix:128} & Retain derivation in appendix \\
Unique binding-cap fertility root & Proved as written on the feasible branch &
\(\vartheta>0\), nonempty stated domain, fixed old-age branch &
\hyperlink{source-appendix}{Appendix:157} & Strictly decreasing FOC establishes
uniqueness \\
Housing wedge and full price of a child & Correct but economically weaker than some
descriptions suggest & Interior fertility; relevant housing KKT condition &
\hyperlink{source-main}{Main:209} & A private shadow-price identity, not a social
resource-cost or inefficiency theorem \\
Cap derivative and sign condition & Conditional & Hold user cost, effective wealth,
tenure, and old-age commitment fixed & \hyperlink{source-appendix}{Appendix:513} & Correct
partial derivative \\
Applying that derivative to the common rental cap & Incomplete; false if treated as the
full derivative & The old rental commitment changes too &
\hyperlink{source-appendix}{Appendix:570} & Add the resource term derived below \\
LTV and liquid-wealth derivatives & Conditional & Only the financing cap binds; prices,
transfers, tenure, and active set fixed & \hyperlink{source-appendix}{Appendix:574} &
Correct local derivatives; no aggregate policy sign \\
Logistic tenure shares & Proved as written & Finite conditional values; positive taste
scale; independent additive taste & \hyperlink{source-main}{Main:130} & Smooths tenure
shares, not every household constraint kink \\
Distribution law, housing clearing, government budget & Conditional & Normalized \(G_t\);
external financial sector; fixed title/occupancy rules &
\hyperlink{source-appendix}{Appendix:218} & Essentially complete; consolidate into one
equilibrium definition \\
Proposition 1: replacement fertility and two-equation steady state & Proved as written &
Positive constant cohort masses and stationary household distribution &
\hyperlink{source-main}{Main:247} & Accounting characterization, not existence \\
Proposition 4: steady-state existence/local uniqueness & Conditional & Fiscal self-map and
uniform contraction; replacement sign change; differentiability and nonsingular Jacobian &
\hyperlink{source-appendix}{Appendix:258} & Correct sufficient theorem; grid checks do not
prove its uniform hypotheses \\
Steady-state price/transfer/population derivatives & Conditional & Regular differentiable
root & \hyperlink{source-appendix}{Appendix:300} & Correct; signs remain unrestricted \\
Gains complementarity and cohort recursion & Proved as written & Nonnegative gains-tax
rate; given demographic law & \hyperlink{source-appendix}{Appendix:329},
\hyperlink{source-main}{Main:285} & Explains persistence; does not guarantee
appreciation \\
Proposition 5: fixed-\(J\) local transition & Conditional & Smooth aggregate branch;
nonsingular \textbf{zero-shock} Jacobian; strictly consistent directional price signs &
\hyperlink{source-appendix}{Appendix:367} & Correct for the homogeneous specialization;
general \(F\) also needs appropriate uniform/integrability regularity \\
Reported numerical transition and horizon stability & Numerically demonstrated but not
proved & Declared specialization and terminal closure &
\hyperlink{source-builder}{Builder:770} & Credible approximation, not an exact infinite
path \\
Young--old marginal-value decomposition & Conditional & Financing-only young wedge;
interior saving and old choices & \hyperlink{source-appendix}{Appendix:399} & Correct
diagnostic \\
A: Proposition 2, compensated reallocation & Conditional & Added bridge finance; targeted
two-date transfers; explicit resale path; real costs counted once &
\hyperlink{source-appendix}{Appendix:421} & Not baseline constrained inefficiency \\
B: constrained buyer outside the feasible tangent set & Proved as written & Fixed
\(P,b,\phi\); rebate unavailable at closing & \hyperlink{source-memo}{Memo:55} & Decisive
restriction on the intended financing interpretation \\
C: zero-cost baseline transition inefficiency & Conditional & Positive-mass uniformly
interior matched subsets; coordinated feasible trades; common current rebate &
\hyperlink{source-memo}{Memo:90} & Valid fiscal-externality result; inapplicable to
displayed constrained buyers \\
C with \(\ell_t>\bar K_t\) & Incomplete; false as a sufficient condition under the stated
instruments & Must specify cost incidence or authorize common financing of costs &
\hyperlink{source-memo}{Memo:125}, \hyperlink{source-ledger}{Ledger:43} & Withdraw the
average-cost sufficiency claim \\
Corollary: model class contains efficient and inefficient allocations & Incomplete as a
self-contained proof & Establish admissible examples; specify continuation resource
accounting and positive-measure strict improvement & \hyperlink{source-memo}{Memo:133} &
Plausible, but contributes little; omit \\
Smooth infinite-horizon proposition & Conditional & \(C^1\) map, invertible next-date
derivative, hyperbolicity, full-rank stable projection & \hyperlink{source-memo}{Memo:195}
& Correct standard local result; does not cover statutory kink \\
Proposed nonsmooth infinite-horizon argument & Incomplete & Precise operator theorem and
uniform switched-system bounds & \hyperlink{source-memo}{Memo:250} & Research outline, not
a completed theorem \\
Taxing intermediaries leaves the decomposition unchanged & Correct but economically weaker
than an invariance claim & Same per-unit tax; distinguish algebraic form from equilibrium
magnitudes & \hyperlink{source-appendix}{Appendix:493} & Identity survives; prices,
tenure, and \(\zeta\) can change \\
\end{longtable}

\endgroup

\end{landscape}\clearpage\subsection*{Details behind the audit}

Three details warrant elaboration.

\textbf{The lifetime reduction is correct, but its branch conditions do real work.}
Writing \(k\) for the old-age saving weight,

\[
(1+k)x+\chi n+uh=\widetilde w,
\]

where

\[
(k,\widetilde w)=
\begin{cases}
(\beta(1+\gamma+\omega_B),w),&\text{interior old choices},\\
(\beta(1+\omega_B),w-q^2r_{t+1}h_R^{\max}),&\text{capped old renter}.
\end{cases}
\]

The old renter correction is implemented correctly. Tenure values must retain the
tenure-dependent constants omitted from this policy reduction; the code correctly
evaluates the original utilities.

\textbf{The common rental-cap derivative needs an additional term.} On the
capped-old-renter branch, let

\[
n_{\widetilde w}
=\frac{\chi}{(1+k)x^2\Delta}>0.
\]

If the same rental cap rises at both ages, then

\[
\frac{dn}{dh_R^{\max}}
=
\left.\frac{\partial n}{\partial\widehat h}\right|_{\widetilde w}
-q^2r_{t+1}n_{\widetilde w}.
\]

The second term reflects additional committed old-age rent. At the final illustrative
steady state, the current-cap partial derivative is approximately \(0.80296\), while the
common-cap derivative is \(0.79348\); an independent finite difference confirms the
latter. This does not overturn the example's sign, but it changes the claimed policy
derivative.

\textbf{The child-unit mapping is unresolved across documents.} The example uses
\(\nu=2\), implying replacement \(n=0.5\). The decision ledger separately says one toy
child unit corresponds to \(2.1\) literal children. Together these imply \(1.05\), not
\(2.1\), literal children per toy household at replacement. The abstract model can use
normalized units, but the biological interpretation and quantitative mapping must be
reconciled. Compare \hyperlink{source-appendix}{Appendix:616} with
\hyperlink{source-ledger}{Ledger:16}.

\section{Welfare verdict}\label{welfare-verdict}

\subsection{A: valid with added credit and
compensation}\label{a-valid-with-added-credit-and-compensation}

The transfer ledger is internally consistent. Define \(u_t^O=qr_t+q\ell_{t+1}\). Then:

\begin{itemize}
\tightlist
\item
  The old incumbent's first-order private loss is zero because \[
  MV_{jt}^{O}-q\tau^pP_t+qP_{t+1}=P_t-\ell_t.
  \]
\item
  The young buyer's net private gain, after purchase, property tax, and eventual taxable
  resale, is \(\zeta_{it}^{O,F}\).
\item
  The government receives incremental revenue \(\ell_t+q\ell_{t+1}\).
\item
  Market-rate lenders break even.
\end{itemize}

With unrestricted implementation of the stated compensation ledger, the available marginal
surplus is therefore

\[
\mathcal S_{ijt}
=\zeta_{it}^{O,F}+\ell_t+q\ell_{t+1}-K_{ijt}.
\]

A strict positive surplus supports a sufficiently small compensated improvement, including
compensation for second-order private losses.

\textbf{But bridge finance is an additional credit technology.} It permits the buyer to
fund the marginal down payment with a pledgeable advance rather than existing liquid
wealth. Its zero present value does not make it institutionally innocuous. A
zero-present-value loan can be extremely valuable precisely because the baseline household
cannot obtain it.

The strongest honest description of A is:

\begin{quote}
Adding pledgeable closing finance and appropriately financed compensation can improve the
allocation of housing among existing households.
\end{quote}

It is not a theorem that the original credit-constrained allocation violates the
efficiency conditions associated with its original feasible set.

\subsection{B: correct, and compatible with A}\label{b-correct-and-compatible-with-a}

At a binding closing constraint,

\[
b_i-(1-\phi)P_th_i=0.
\]

Holding \(b_i,P_t,\phi\) fixed, feasibility of a one-sided perturbation requires

\[
-(1-\phi)P_t\dot h_i\geq0,
\qquad\text{hence}\qquad \dot h_i\leq0.
\]

The proposed recipient cannot receive more owned housing under that institution.

This does \textbf{not} establish constrained efficiency of the entire economy. Other
reallocations or equilibrium price changes could matter. If price changes,

\[
P_t\dot h_i+h_i\dot P_t\leq0,
\]

so falling prices can expand affordable housing. B rules out the particular fixed-price
transfer, not all possible policies.

\subsection{C: the zero-cost result works}\label{c-the-zero-cost-result-works}

The proof can be made concrete without targeted transfers.

Consider an unconstrained young buyer receiving \(\epsilon\) extra housing. Hold fertility
and subsequent retained housing fixed. Before receiving the incremental common rebate,
choose:

\[
\Delta a'
=\bigl[\phi P_t-q(P_{t+1}-\ell_{t+1})\bigr]\epsilon,
\qquad
\Delta c=-u_t^O\epsilon.
\]

These changes satisfy the young budget and leave old-age resources unchanged before the
additional future rebate. Slack closing and saving constraints make them feasible. The
young household's first-order utility change is zero because \(MV^Y=u_t^O\).

For the old seller, reduce retention by \(\epsilon\), hold total estate expenditure fixed,
and increase current consumption by

\[
\Delta c^2=(qr_t-\ell_t)\epsilon.
\]

The old household's first-order private utility change is also zero.

Now perform these transfers within matched subsets of mass \(m\). Current gains-tax
revenue increases by \(m\ell_t\epsilon\), generating the common rebate

\[
\Delta T_t=\frac{m\ell_t}{Y_t+O_t}\epsilon.
\]

Private losses on the matched subsets are uniformly \(O(\epsilon^2)\); the rebate benefit
is \(O(\epsilon)\). Thus, for sufficiently small \(\epsilon\), it dominates those losses.

This also addresses the omitted continuation accounting:

\begingroup\small\setlength{\tabcolsep}{5pt}\renewcommand{\arraystretch}{1.17}

\begin{longtable}[]{@{}
  >{\raggedright\arraybackslash}p{(\columnwidth - 2\tabcolsep) * \real{0.2600}}
  >{\raggedright\arraybackslash}p{(\columnwidth - 2\tabcolsep) * \real{0.7400}}@{}}
\toprule\noalign{}
\begin{minipage}[b]{\linewidth}\raggedright
Concern
\end{minipage} & \begin{minipage}[b]{\linewidth}\raggedright
Resolution in the zero-cost construction
\end{minipage} \\
\midrule\noalign{}
\endhead
\bottomrule\noalign{}
\endlastfoot
Buyer's future resale & The extra unit is sold when the buyer becomes old; baseline
retained housing is unchanged \\
Tax basis & Extra housing has acquisition basis \(P_t\) \\
Future tax & Generates \(m\ell_{t+1}\epsilon\) at \(t+1\) \\
Baseline estate sale & The original incumbent would have sold that unit through the estate
without gains tax \\
Property tax & Aggregate title/occupancy changes leave the total property-tax base
unchanged \\
Old incumbent's death & Current appreciation supplies compensation while the incumbent is
alive \\
Other current households & Can consume the additional common current rebate \\
Future households & Can receive the additional \(t+1\) rebate while retaining baseline
choices; no loss is required \\
Prices & The constructed allocation preserves housing clearing at the original prices; it
is not a newly decentralized policy equilibrium \\
\end{longtable}

\endgroup

\textbf{The strongest baseline theorem is therefore:}

\begin{quote}
At fixed fertility, population, and prices, current realization taxation creates a local
Pareto-improving feasible allocation when a positive mass of sufficiently interior buyers
and sellers can make costless coordinated trades and the resulting current revenue is
rebated equally.
\end{quote}

Mathematically, the relevant sufficient condition in this construction is

\[
\ell_t>0,\qquad K=0,
\]

together with the stated positive-mass and uniform-interiority assumptions.

It is a \textbf{realization-tax allocative distortion with a fiscal externality}. It is
not merely a pecuniary externality: prices are held fixed. Nor is it a costless wash-sale
trick: housing services actually move between households with different marginal
valuations. However, it does not depend on fertility and does not establish that
financially constrained families should receive the housing.

It also requires coordination of trades. The common rebate alone does not decentralize the
proposed allocation as individually optimal behavior.

\subsection{The average-cost extension does not
follow}\label{the-average-cost-extension-does-not-follow}

The condition \(\ell_t>\bar K_t\) identifies an aggregate surplus. It does not identify
who pays the costs or ensure compensation under a common rebate.

For example, suppose \(m/(Y_t+O_t)=0.1\), current tax revenue per transferred unit is
\(1\), future gains are zero, and each buyer incurs a transaction cost \(0.5\). Then

\[
\ell_t=1>\bar K_t=0.5,
\]

but each buyer receives only \(0.1\epsilon\) through the current rebate and incurs
\(0.5\epsilon\) in cost. The first-order loss is not repaired.

The extension could work if the government paid the transaction costs before distributing
a net common rebate. But that introduces an expenditure instrument absent from the stated
government budget. Alternatively, participant-specific cost bounds could establish that
each party's rebate covers its own cost. An average-cost inequality cannot substitute for
those conditions.

The memo's wording is incomplete; the decision ledger's stronger reading of it as
sufficiency is incorrect.

\subsection{What should remain in the paper?}\label{what-should-remain-in-the-paper}

\textbf{Only the marginal-value and child-cost decompositions. No welfare proposition in
the main text.}

C does not apply to the displayed numerical construction: all young owners face binding
financing limits, and young renters also face binding rental limits. A applies only after
adding institutions that are not part of the quantitative exercise.

For policy welfare, a more natural object is the \textbf{distribution of
consumption-equivalent changes for households already alive}, reported by age, wealth, and
tenure. This permits transparent reporting of winners and losers without inventing
transfers that the policy does not provide.

Comparisons across different future populations require an explicit additional criterion.
They do not necessarily require arbitrary utilitarian weights, but they do require
decisions about potential people and their identities. That distinction is central in
\href{https://onlinelibrary.wiley.com/doi/abs/10.1111/j.1468-0262.2007.00781.x}{Golosov,
Jones, and Tertilt (2007)} and
\href{https://onlinelibrary.wiley.com/doi/full/10.3982/TE2138}{Pérez-Nievas, Conde-Ruiz,
and Giménez (2019)}.

If the desired result is specifically an improvement for constrained young buyers, the
smallest algebraic addition is a pledgeable down-payment advance. Its economic
justification must explain why the relevant authority can supply credit that private
lenders cannot. Otherwise the theorem simply restates that relaxing a binding constraint
is valuable. I would prefer evaluating an actual financing policy with its costs and
distributional effects over adding a friction solely to obtain a Pareto theorem.

\section{Transition verdict}\label{transition-verdict}

The fixed-\(J\) proposition proves a narrow but legitimate result:

\begin{enumerate}
\def\labelenumi{\arabic{enumi}.}
\tightlist
\item
  Select a smooth appreciation branch.
\item
  Assume a nonsingular stacked Jacobian at the zero-shock steady state.
\item
  Verify that the resulting directional price changes have the selected signs.
\item
  Obtain a locally unique solution \textbf{on that branch}, for a sufficiently small shock
  and that fixed horizon.
\end{enumerate}

It does not establish another branch's nonexistence, a shock-size neighborhood uniform in
\(J\), or convergence as \(J\to\infty\).

The numerical evidence is strong within its proper scope:

\begingroup\small\setlength{\tabcolsep}{5pt}\renewcommand{\arraystretch}{1.17}

\begin{longtable}[]{@{}
  >{\raggedright\arraybackslash}p{(\columnwidth - 2\tabcolsep) * \real{0.6400}}
  >{\raggedright\arraybackslash}p{(\columnwidth - 2\tabcolsep) * \real{0.3600}}@{}}
\toprule\noalign{}
\begin{minipage}[b]{\linewidth}\raggedright
Check
\end{minipage} & \begin{minipage}[b]{\linewidth}\raggedright
Result
\end{minipage} \\
\midrule\noalign{}
\endhead
\bottomrule\noalign{}
\endlastfoot
Existing endpoint tests & All four pass \\
Independent original-program optimizations & All 12 converge; largest objective
discrepancy \(3.73\times10^{-14}\) \\
Initial/final illustrative prices & \(0.342683563,\ 0.484405426\) \\
Initial/final young owner shares & \(0.748566127,\ 0.514929506\) \\
Average fertility at both endpoints & \(0.5\) \\
Horizon-28 maximum scaled residual & Approximately \(6.50\times10^{-10}\) \\
Fresh horizon-24/28 difference, first nine dates & Approximately \(2.22\times10^{-15}\) \\
Saved maximum reported terminal-state gap & Approximately \(1.87\times10^{-9}\) \\
\end{longtable}

\endgroup

The maintained active sets pass. In particular, old-owner retention and estate constraints
remain slack, and the old-renter cap remains binding.

There are three limits to this evidence:

\begin{itemize}
\tightlist
\item
  \textbf{The reported branch Jacobian is evaluated at the finite-shock solution.} It does
  not verify the zero-shock directional theorem. See
  \hyperlink{source-builder}{Builder:986}.
\item
  \textbf{The horizon comparator checks selected aggregates over dates 0--8.} The
  terminal-state summary also omits the acquisition-price coordinate. See
  \hyperlink{source-builder}{Builder:1397}.
\item
  \textbf{Tiny tail appreciation signs are not robustly certified.} The smallest price
  change is around \(1.84\times10^{-10}\). Dates 0--3 and 7--8 are the dates with gains
  taxes above \(10^{-4}\); they are not the exhaustive list of positive floating-point
  gains.
\end{itemize}

The generalized-root calculations also reproduce: seven stable roots, two finite unstable
roots, and one infinite root associated with the static transfer. Eliminating that
transfer removes the infinite root. The final-state sampled stable modulus is at most
about \(0.547\).

That supports the frozen linearizations. It does not prove stability when the economy
switches between tax regimes. Stable or hyperbolic individual matrices do not supply the
necessary bounds on their products.

\textbf{An infinite-horizon theorem is not economically necessary for the proposed
main-text role.} A clearly labeled numerical transition can demonstrate demographic
propagation and several episodes of appreciation.

If the paper insists on an exact existence claim, the weakest useful target is existence
of a convergent path for the homogeneous specialization, without global uniqueness. A
limit argument would need actual solutions for arbitrarily large horizons, compactness and
continuity sufficient to pass equilibrium conditions to the limit, and a uniform tail
condition selecting the final steady state. The existing horizon checks do not establish
these assumptions.

A local existence-and-uniqueness result for the exact kinked model requires substantially
stronger control of the admissible switching dynamics and the initial-state boundary
condition. The memo correctly identifies the difficulty but has not supplied the theorem.

One nuance: smoothing the gains schedule need not remove realization-based lock-in
altogether. A smooth positive tax on gains combined with stepped-up basis still rewards
retention. It removes the kink; setting \(\tau^g=0\) removes the gains-tax mechanism.
Those are different changes.

\section{Economic assessment}\label{economic-assessment}

\textbf{The useful economics is the distinction between obtaining ownership and obtaining
child-relevant housing services.}

The identity

\[
\frac{\vartheta x}{n}
=\chi+\kappa(u+\zeta)
\]

makes that distinction transparent. Children use goods and space; housing constraints can
raise the private shadow cost of the latter. Relaxing a housing cap can raise fertility,
but paying for additional housing can offset the space benefit.

Log utility is doing useful analytical work: it gives simple expenditure shares, a unique
conditional solution, and a readable comparative static. But it also imposes strong
behavior:

\begin{itemize}
\tightlist
\item
  Everyone chooses strictly positive continuous fertility.
\item
  There is no childlessness margin, birth timing, or sequential parity decision.
\item
  Conditional unconstrained fertility is proportional to effective wealth.
\item
  Housing requirements are technologically imposed rather than chosen through child
  quality or household production.
\end{itemize}

These limitations are acceptable for a mechanism model. They preclude presenting its
elasticities as structural counterparts of the full sequential model.

Ownership adds something essential only if the question concerns \textbf{liquidity, resale
wealth, or title-dependent taxes}. A renter-only model can already establish the
goods-versus-space fertility mechanism. Ownership earns its place through the financing
constraint; gains taxation requires a separate justification.

The fixed housing stock is a reasonable short-run simplification. It becomes restrictive
as a long-run theory. Since the two steady-state equations determine \(P,T\) independently
of \(\bar H\), increasing the stock simply scales steady-state population in this
specification. The result is clean, but much of it comes from replication of identical
economic environments.

The replacement-fertility result is useful as a diagnostic restriction, not a substantive
explanation of fertility:

\begin{quote}
A positive constant population in a closed reproduction system must replace itself.
\end{quote}

That conclusion does not depend on logarithmic utility or fixed housing supply. Housing
determines whether a supporting equilibrium exists and what population scale it supports.
It does not guarantee that prices can always restore replacement.

The two-period structure clarifies the distinction between a living sale and an estate
sale. It also makes tax lock-in unusually stark: death is certain, only one retention
period remains, household landlords are prohibited, and tax bases have only one
acquisition vintage. This is a useful illustration of avoidance incentives, not a
convincing account of realistic housing turnover by itself.

The common rebate is a transparent fiscal closure. It is also economically powerful. In C,
it supplies the first-order benefit that repairs the participants' second-order losses. It
should therefore be presented as a maintained institution, not neutral bookkeeping.

The ownership taste is a reasonable smoothing device. One additive draw is sufficient;
richer taste machinery would add little. Nevertheless, in the homogeneous numerical
example, the taste distribution is doing substantive aggregate work: mixed tenure is not
generated by heterogeneous incomes or wealth.

A permanent preference shock is a clean experiment. \textbf{Its upward direction is
selected to produce appreciation.} The main quantitative application concerns a fertility
decline, which need not generate the same gains-tax mechanism. That asymmetry should
prevent the tax example from becoming the paper's general account of demographic change.

The two proposed theory pieces share an environment but currently provide \textbf{two
distinct lessons}:

\begin{itemize}
\tightlist
\item
  A conditional statement about housing access and private valuations.
\item
  A numerical statement about demographic propagation and transitional tax bases.
\end{itemize}

They become a coherent paper contribution only if the tax mechanism matters for the
empirical or quantitative exercise. At present, it does not.

The quantitative connections are narrower and stronger:

\begin{itemize}
\tightlist
\item
  Housing services, wealth, and budgets affect fertility.
\item
  Tenure changes alone need not change fertility materially.
\item
  Population evolves through cohort propagation.
\item
  The existence of a positive terminal population must be checked, not presumed.
\end{itemize}

The toy model cannot establish a capital-gains mechanism absent from the quantitative code
or repair its unverified positive closed endpoint. The live quantitative distinction
between housing-service responses and financing-only tenure responses is documented in
\hyperlink{source-status}{Status:90}.

My assessment is therefore:

\begingroup\small\setlength{\tabcolsep}{5pt}\renewcommand{\arraystretch}{1.17}

\begin{longtable}[]{@{}
  >{\raggedright\arraybackslash}p{(\columnwidth - 2\tabcolsep) * \real{0.3900}}
  >{\raggedright\arraybackslash}p{(\columnwidth - 2\tabcolsep) * \real{0.6100}}@{}}
\toprule\noalign{}
\begin{minipage}[b]{\linewidth}\raggedright
Standard
\end{minipage} & \begin{minipage}[b]{\linewidth}\raggedright
Assessment
\end{minipage} \\
\midrule\noalign{}
\endhead
\bottomrule\noalign{}
\endlastfoot
Intellectually correct & Most core results, with the qualifications above \\
Useful for the paper & Child-cost mechanism and demographic accounting \\
Worth five main-text pages as currently assembled & No \\
Worth presenting in a seminar & One mechanism slide; possibly one demographic diagram \\
Worth an appendix & Household derivations and numerical transition \\
Worth an online supplement & Gains-tax welfare extension and advanced transition
diagnostics \\
\end{longtable}

\endgroup

\section{Taste and exposition}\label{taste-and-exposition}

The current rendered main section is \textbf{five pages}, followed by nine appendix pages
and a references page. The older review's six-page description is stale. I inspected the
current package and formalization memo: \hyperlink{source-package-pdf}{Theory PDF}
\hyperlink{source-memo-pdf}{Memo PDF}

The writing is substantially clearer than a typical overformalized model note. The
weaknesses are mainly allocation of attention and claim hierarchy.

\begin{enumerate}
\def\labelenumi{\arabic{enumi}.}
\item
  \textbf{The opening promises more than the section establishes.} ``Move the economy
  between positive steady states'' precedes the later admission that no infinite-horizon
  result is established. Start with the mechanism and describe the transition as an
  illustration. See \hyperlink{source-main}{Main:6}.
\item
  \textbf{Some primitives appear too late.} The ownership taste, discount factor, rebate
  rule, and closing chronology belong in the environment. They currently arrive inside or
  after the household problems.
\item
  \textbf{The rental intermediary is explained well.} Keep its short cash-flow account.
  The sentence about taxing intermediaries can move to an appendix footnote.
\item
  \textbf{Stepped-up basis is understandable, but needs one ordinary example.} A house
  bought for 100 and worth 150 generates taxable gains of 50 on a living sale; an
  immediate estate sale after resetting basis to 150 generates zero. State that this is
  the model's tax rule.
\item
  \textbf{The old-owner budget is too compressed for its importance.} Show sale proceeds,
  retained housing, and financial estate once. That would make the estate floor
  intelligible without requiring a reader to reverse-engineer it.
\item
  \textbf{The order of household problems should make continuation values available before
  they are used.} A compact old-age problem followed by the young problem is the cleanest
  recursive presentation.
\item
  \textbf{Equilibrium should be defined once.} Housing clearing and fiscal balance appear
  twice within a short span. Replace the repetition with one numbered definition. See
  \hyperlink{source-main}{Main:234}.
\item
  \textbf{The notation is mostly defensible.} Keep \(\chi,\kappa,P,r,T\). Use a visibly
  distinct symbol for old households rather than overloading \(O\) for old age and
  ownership. ``Net-of-child goods'' and ``net-of-child space'' are clearer than repeatedly
  saying ``usable bundles.'' Explain \(\zeta\) as the extra willingness to pay per unit of
  space; it is not itself the value of one dollar of liquidity.
\item
  \textbf{The most useful comparative static is buried.} The space benefit versus
  goods-cost tradeoff belongs in the main mechanism discussion. The reduced maximization
  problem and expenditure-share catalogue can move out.
\item
  \textbf{The final subsection's title promises fertility but then fixes fertility.}
  ``Intergenerational reallocation and fertility'' is misleading for that comparison. See
  \hyperlink{source-main}{Main:305}.
\item
  \textbf{Qualifications currently carry too much of the argument.} This is not solved by
  deleting caveats. Choose a narrower headline result so fewer caveats are needed.
\item
  \textbf{The appendix is too long relative to the contribution.} Its length comes from
  retaining several different ambitions. The household derivations are worth keeping; the
  contraction theorem, branchwise transition theorem, welfare ledger, and saddle apparatus
  should not all receive equal prominence.
\end{enumerate}

The prose occasionally sounds procedural: ``isolates two implications,'' repeated
statements that an appendix establishes something, and the succession of generic theorem
announcements. I would not infer authorship from this. The editorial problem is that the
text describes the machinery more often than it explains why the reader needs the result.

\textbf{Figures.} The transition figure is informative at full size but too small at its
current half-text-width placement. Its four panels become miniature plots. If retained,
give it full width in the appendix, identify the pre-shock point, label the owner share as
the young owner share, and explain the illustrative fertility units.

The wedge figure is readable and accurately displays a valuation gap. It is largely
redundant with one equation, and its long labels dominate the graphic. I would omit it
from the paper. It should never be used as visual evidence of a measured welfare loss.

Neither figure directly shows the fertility response to housing access---the analytical
result most relevant to the paper.

\section{Minimum viable theory}\label{minimum-viable-theory}

I would retain the following ingredients:

\begin{itemize}
\tightlist
\item
  Two overlapping adult cohorts and explicit cohort accounting.
\item
  Goods and housing requirements per child.
\item
  Saving and a common borrowing/lending return.
\item
  Renting versus ownership, with a clearly stated financing restriction.
\item
  A simple rental-size restriction, if that institution is part of the intended mechanism.
\item
  Old-age housing and estate motives sufficient to describe lifecycle demand.
\item
  A common property-tax rebate if property-tax policy remains central.
\item
  Fixed housing for the short-run illustration.
\item
  No capital-gains tax in the main benchmark.
\item
  One ownership taste distribution only if mixed tenure is needed for the illustration.
\end{itemize}

\textbf{Five-page main-text ceiling}

\begingroup\small\setlength{\tabcolsep}{5pt}\renewcommand{\arraystretch}{1.17}

\begin{longtable}[]{@{}
  >{\raggedright\arraybackslash}p{(\columnwidth - 2\tabcolsep) * \real{0.2800}}
  >{\raggedright\arraybackslash}p{(\columnwidth - 2\tabcolsep) * \real{0.7200}}@{}}
\toprule\noalign{}
\begin{minipage}[b]{\linewidth}\raggedright
Space
\end{minipage} & \begin{minipage}[b]{\linewidth}\raggedright
Retained content
\end{minipage} \\
\midrule\noalign{}
\endhead
\bottomrule\noalign{}
\endlastfoot
About one page & Environment, timing, and institutions \\
About one page & Compact household problems and one equilibrium definition \\
About one page & Child-cost identity and conditional housing-cap response \\
About one page & Cohort propagation, replacement accounting, and the possibility that no
positive steady state exists \\
Remaining space & Interpretation and, only if needed, one clearly labeled numerical
illustration \\
\end{longtable}

\endgroup

The retained theorem set would be:

\begin{enumerate}
\def\labelenumi{\arabic{enumi}.}
\tightlist
\item
  \textbf{One household proposition:} the full private price of a child and the
  conditional sign of the housing-access response.
\item
  \textbf{One demographic lemma:} positive closed stationarity requires \(\nu\bar n=1\);
  housing clearing determines scale.
\item
  \textbf{No welfare proposition.}
\item
  \textbf{No infinite-horizon transition theorem.}
\end{enumerate}

A useful nonexistence statement is simply:

\[
\sup_{P\ \text{on the feasible fiscal locus}}
\nu\bar n(P,\mathcal T(P);\vartheta)<1
\quad\Longrightarrow\quad
\text{no positive closed steady state}.
\]

It is modest, but directly disciplines interpretation of the quantitative model.

\textbf{Appendix}

Retain the original-budget derivations, branch conditions, corrected cap derivatives,
aggregation, and one reproducible numerical packet. State the finite terminal closure
precisely. Compress generic existence and implicit-function arguments.

\textbf{Park outside the paper}

The expanded-credit Pareto theorem, the baseline tax/rebate theorem, the nonsmooth
operator argument, the full generalized-root discussion, and the wedge graphic. Preserve
them as research material; do not give them paper space merely because they are
technically available.

\section{Priority ladder}\label{priority-ladder}

\begingroup\small\setlength{\tabcolsep}{5pt}\renewcommand{\arraystretch}{1.17}

\begin{longtable}[]{@{}
  >{\raggedright\arraybackslash}p{(\columnwidth - 4\tabcolsep) * \real{0.1800}}
  >{\raggedright\arraybackslash}p{(\columnwidth - 4\tabcolsep) * \real{0.6000}}
  >{\raggedright\arraybackslash}p{(\columnwidth - 4\tabcolsep) * \real{0.2200}}@{}}
\toprule\noalign{}
\begin{minipage}[b]{\linewidth}\raggedright
Priority
\end{minipage} & \begin{minipage}[b]{\linewidth}\raggedright
Next concrete action
\end{minipage} & \begin{minipage}[b]{\linewidth}\raggedright
Author decision?
\end{minipage} \\
\midrule\noalign{}
\endhead
\bottomrule\noalign{}
\endlastfoot
\textbf{Fatal before use} & Remove any claim that the financing wedge proves baseline
constrained inefficiency; distinguish A's expanded credit institution from B's baseline
feasible set & \textbf{Yes:} choose intended welfare benchmark \\
\textbf{Fatal before use} & Withdraw \(\ell_t>\bar K_t\) as sufficient under common
rebates unless cost financing/incidence is explicitly supplied & Only if retaining a
positive-cost welfare theorem \\
\textbf{Fatal before use} & Remove literal claims that the computation proves an exact
infinite-horizon transition or validates the quantitative closed endpoint & No; this is a
claim-boundary correction \\
\textbf{Important before circulating} & Decide whether capital-gains taxation belongs in
the main analytical model & \textbf{Yes} \\
\textbf{Important before circulating} & Distinguish a young-only cap relaxation from
changing the rental cap at both ages; add the omitted resource term & No \\
\textbf{Important before circulating} & Reconcile \(\nu\), normalized fertility units, and
the claimed mapping to literal children & \textbf{Yes:} choose the normalization \\
\textbf{Important before circulating} & Consolidate equilibrium and state the external
financial sector, timing, and title/occupancy restrictions explicitly & No, unless
changing institutions \\
\textbf{Important before circulating} & Describe the finite-shock Jacobian honestly;
either verify zero-shock hypotheses or leave the proposition explicitly conditional &
No \\
\textbf{Useful improvement} & Include acquisition basis in terminal diagnostics and attach
a threshold to ``positive gains'' dates & No \\
\textbf{Useful improvement} & Move the informative fertility comparative static into the
main discussion; shorten the branch catalogue there & No substantive decision \\
\textbf{Useful improvement} & Enlarge or demote the transition figure; omit the redundant
wedge figure & Editorial decision \\
\textbf{Optional} & Add selected tests for unused corners if the code is to become a
general solver & No need for the maintained example \\
\textbf{Optional} & Evaluate current-household consumption-equivalent policy changes under
actual fiscal and lending rules & \textbf{Yes:} choose policy and welfare reporting
scope \\
\textbf{Optional} & Pursue a certified infinite-horizon result & \textbf{Yes; my
recommendation is not to do so} \\
\end{longtable}

\endgroup

\section{Final recommendation}\label{final-recommendation}

If this were my job-market paper, I would make the analytical section explain one
proposition: \textbf{housing affects fertility through the services and resources
available to families; ownership matters insofar as it changes access to those services or
resources.} That proposition connects directly to the quantitative finding that financing
policy can move tenure much more than fertility.

I would retain demographic accounting to prevent an unsupported long-run interpretation. I
would not use replacement fertility as a major theoretical result, and I would not let a
successful toy transition imply a terminal equilibrium that the quantitative model has not
established.

I would remove the main-text Pareto proposition. A is a valid theorem about additional
credit and compensation; C is a valid zero-cost fiscal-externality construction for
unconstrained recipients. Neither establishes the intended claim about the displayed
constrained young families under the baseline institutions.

The analytical model is worth keeping only after that narrowing. Its core is useful. Its
present collection of welfare, tax-basis, steady-state, and nonsmooth-transition ambitions
asks the reader to learn substantially more machinery than the paper presently uses.

\clearpage
\appendix
\section{Source guide}
Source references in the report link to the entries below. Line numbers identify the versions reviewed; subsequent edits may move them. All paths are absolute.

\Needspace{5\baselineskip}
\hypertarget{source-main}{\textbf{Main -- Main analytical section}}\par
{\small\nolinkurl{/Users/tommasodesanto/Desktop/Projects/Fertility/Fertility_Spring26/latex/simplified_olg_paper_theory_section.tex}}\par
\vspace{0.35em}

\Needspace{5\baselineskip}
\hypertarget{source-appendix}{\textbf{Appendix -- Proof and computational appendix}}\par
{\small\nolinkurl{/Users/tommasodesanto/Desktop/Projects/Fertility/Fertility_Spring26/latex/simplified_olg_paper_theory_appendix.tex}}\par
\vspace{0.35em}

\Needspace{5\baselineskip}
\hypertarget{source-builder}{\textbf{Builder -- Numerical construction}}\par
{\small\nolinkurl{/Users/tommasodesanto/Desktop/Projects/Fertility/Fertility_Spring26/code/model/tools/build_simplified_olg_mixed_tenure_theory.py}}\par
\vspace{0.35em}

\Needspace{5\baselineskip}
\hypertarget{source-memo}{\textbf{Memo -- Competing formalization memo}}\par
{\small\nolinkurl{/Users/tommasodesanto/Desktop/Projects/Fertility/Fertility_Spring26/latex/JMP_DS_suggestions/simplified_olg_formalization.tex}}\par
\vspace{0.35em}

\Needspace{5\baselineskip}
\hypertarget{source-ledger}{\textbf{Ledger -- Active decision ledger}}\par
{\small\nolinkurl{/Users/tommasodesanto/Desktop/Projects/Fertility/Fertility_Spring26/docs/model/ACTIVE_DECISION_LEDGER.md}}\par
\vspace{0.35em}

\Needspace{5\baselineskip}
\hypertarget{source-status}{\textbf{Status -- Canonical quantitative status}}\par
{\small\nolinkurl{/Users/tommasodesanto/Desktop/Projects/Fertility/Fertility_Spring26/CALIBRATION_STATUS.md}}\par
\vspace{0.35em}

\Needspace{5\baselineskip}
\hypertarget{source-package-pdf}{\textbf{Theory PDF -- Rendered analytical package}}\par
{\small\nolinkurl{/Users/tommasodesanto/Desktop/Projects/Fertility/Fertility_Spring26/latex/simplified_olg_paper_theory_package.pdf}}\par
\vspace{0.35em}

\Needspace{5\baselineskip}
\hypertarget{source-memo-pdf}{\textbf{Memo PDF -- Rendered formalization memo}}\par
{\small\nolinkurl{/Users/tommasodesanto/Desktop/Projects/Fertility/Fertility_Spring26/output/pdf/simplified_olg_formalization.pdf}}\par
\vspace{0.35em}

\subsection*{Research references}
Golosov, Mikhail, Larry E. Jones, and Mich\`ele Tertilt (2007). ``Efficiency with Endogenous Population Growth.'' \emph{Econometrica} 75(4), 1039--1071. \href{https://doi.org/10.1111/j.1468-0262.2007.00781.x}{doi:10.1111/j.1468-0262.2007.00781.x}.

P\'erez-Nievas, Mikel, Jos\'e I. Conde-Ruiz, and Eduardo L. Gim\'enez (2019). ``Efficiency and endogenous fertility.'' \emph{Theoretical Economics} 14(2), 475--512. \href{https://doi.org/10.3982/TE2138}{doi:10.3982/TE2138}.

\end{document}
